\documentclass{ctexart}

% preamble.tex
\ctexset{
    section = {
      format = {\Large\bfseries\raggedright},
      indent = {0pt}
     }
}

\usepackage{microtype}
\usepackage{lmodern}
\usepackage{enumitem}
\usepackage[a4paper, margin=1in]{geometry}
\usepackage{xcolor}
\usepackage[cache=true]{minted}
\usepackage[most]{tcolorbox}
\usepackage{tikz}
\tcbuselibrary{minted, skins, breakable}

\setminted{
    fontsize=\small,
    style=default,
    linenos=false,
    breaklines=true,
    fontfamily=tt
}

\newtcolorbox{problembox}[2]{
    enhanced,
    breakable,
    lines before break=3,
    sharp corners,
    boxrule=1.5pt,
    colframe=orange!90!black,
    colback=white,
    colbacktitle=white,
    coltitle=black,
    fonttitle=\bfseries,
    title={Problem (#1): #2},
    titlerule=0pt,
    overlay={
            \draw[black, line width=0.4pt] ([xshift=15pt, yshift=3pt]title.south west) -- ([xshift=-15pt, yshift=3pt]title.south east);
        },
    toptitle=6pt,
    bottomtitle=6pt,
    top=8pt,
    left=10pt,
    right=10pt,
    bottom=10pt
}

\newtcolorbox{tipbox}[1]{
    enhanced,
    breakable,
    lines before break=3,
    sharp corners,
    boxrule=2pt,               % 边框稍微加粗一点以匹配截图的视觉效果
    colframe=blue,             % 边框颜色改为纯蓝色
    colback=white,
    colbacktitle=white,
    coltitle=black,
    fonttitle=\bfseries,
    title={Low-Resource/Downscaling Tip: #1}, % 固定前缀，#1 接收自定义标题
    titlerule=0pt,
    overlay={
            % 调整了 xshift=10pt，使其与下方正文的 left=10pt 和 right=10pt 完美对齐
            \draw[black, line width=0.4pt] ([xshift=10pt, yshift=2pt]title.south west) -- ([xshift=-10pt, yshift=2pt]title.south east);
        },
    toptitle=8pt,
    bottomtitle=6pt,
    top=8pt,
    left=10pt,
    right=10pt,
    bottom=10pt
}

\newtcolorbox{ideabox}[1]{
    enhanced,
    breakable,
    lines before break=3,
    sharp corners,
    boxrule=3pt,
    colframe=yellow,
    colback=white,
    colbacktitle=white,
    coltitle=black,
    fonttitle=\bfseries,
    title={Insights: #1}, % 固定前缀，#1 接收自定义标题
    titlerule=0pt,
    overlay={
            % 调整了 xshift=10pt，使其与下方正文的 left=10pt 和 right=10pt 完美对齐
            \draw[black, line width=0.4pt] ([xshift=10pt, yshift=2pt]title.south west) -- ([xshift=-10pt, yshift=2pt]title.south east);
        },
    toptitle=8pt,
    bottomtitle=6pt,
    top=8pt,
    left=10pt,
    right=10pt,
    bottom=10pt
}

\newtcolorbox{answerbox}{
    enhanced,
    breakable,
    frame hidden,
    colback=blue!3,
    borderline west={2pt}{0pt}{blue!60!black},
    top=8pt,
    bottom=8pt,
    left=12pt,
    right=10pt,
    before skip=1em,
    after skip=1.5em,
    after upper=\vspace*{0.8em}   % <--- 【终极绝招】星号 \vspace* 强制保留这部分留白，绝不会被列表吃掉！
}

\begin{document}

\section{Assignment Overview}

\paragraph{What you will implement}

\begin{enumerate}[itemsep=0pt, parsep=0pt, topsep=0pt]
    \item Byte-pair encoding (BPE) tokenizer \hfill \tcbox[on line, boxsep=0pt, left=2pt, right=2pt, top=2pt, bottom=2pt, colback=green!10, colframe=green!50!black, fontupper=\tiny\bfseries]{DONE}
    \item Transformer language model (LM)
    \item The cross-entropy loss function and the AdamW optimizer
    \item The training loop, with support for serializing and loading model and optimizer state
\end{enumerate}

\paragraph{What you will run}

\begin{enumerate}[itemsep=0pt, parsep=0pt, topsep=0pt]
    \item Train a BPE tokenizer on the TinyStories dataset. \hfill \tcbox[on line, boxsep=0pt, left=2pt, right=2pt, top=2pt, bottom=2pt, colback=green!10, colframe=green!50!black, fontupper=\tiny\bfseries]{DONE}
    \item Run your trained tokenizer on the dataset to convert it into a sequence of integer IDs.
    \item Train a Transformer LM on the TinyStories dataset.
    \item Generate samples and evaluate perplexity using the trained Transformer LM.
    \item Train models on OpenWebText and submit your attained perplexities to a leaderboard.
\end{enumerate}


\newpage

\section{Byte-Pair Encoding (BPE) Tokenizer}
\subsection{The Unicode Standard}

In Python, you can use the \mintinline{python}{ord()} function to convert a single Unicode character into its integer representation.
The \mintinline{python}{chr()} function converts an integer Unicode code point into a string with the corresponding character.

\begin{problembox}{unicode1}{Understanding Unicode (1 point)}

    \begin{enumerate}[label=(\alph*), itemsep=0.5em, topsep=0pt]

        \item What Unicode character does \texttt{chr(0)} return? \\
              \textbf{Answer:} \mintinline{python}{'\x00'}

        \item How does this character's string representation \mintinline{python}{__repr__()} differ from its printed representation? \\
              \textbf{Answer:} The \mintinline{python}{__repr__()} is designed to be \textbf{unambiguous} and is mainly for developers, showing the technical details of the object. On the other hand, the printed representation is meant to be \textbf{readable} and user-friendly.

        \item What happens when this character occurs in text? It may be helpful to play around with the following in your Python interpreter and see if it matches your expectations:

              \textbf{Answer:} Code below.

              \begin{minted}[xleftmargin=1.5em, framesep=5pt]{pycon}
>>> chr(0)
'\x00'
>>> print(chr(0))

>>> "this is a test" + chr(0) + "string"
'this is a test\x00string'
>>> print("this is a test" + chr(0) + "string")
this is a teststring
\end{minted}

    \end{enumerate}

\end{problembox}

\newpage
\subsection{Unicode Encodings}

\begin{itemize}[itemsep=0pt, parsep=0pt, topsep=0pt]
    \item UTF-8: The Efficient Variable-Length Encoder, It uses 1 to 4 bytes per character.
    \item UTF-16: The Middle Ground, It uses either 2 or 4 bytes per character.
    \item UTF-32: The Simple Fixed-Length Encoder, It uses exactly 4 bytes for every character, regardless of what it is.
\end{itemize}

\begin{problembox}{unicode2}{Unicode Encodings (3 point)}

    \begin{enumerate}[label=(\alph*), itemsep=0.5em, topsep=0pt]

        \item What are some reasons to prefer training our tokenizer on UTF-8 encoded bytes, rather than UTF-16 or UTF-32? It may be helpful to compare the output of these encodings for various input strings.\\
              \textbf{Answer:} \mintinline{python}{'\x00'}

        \item Consider the following (incorrect) function, which is intended to decode a UTF-8 byte string into a Unicode string. Why is this function incorrect? Provide an example of an input byte string that yields incorrect results. \\
              \vspace{-1.5em}
              \begin{minted}[xleftmargin=1.5em, framesep=5pt]{python}
def decode_utf8_bytes_to_str_wrong(bytestring: bytes): 
    return "".join([bytes([b]).decode("utf-8") for b in bytestring])
\end{minted}

              \vspace{-1em}

              \begin{minted}[xleftmargin=1.5em, framesep=5pt]{pycon}
>>> decode_utf8_bytes_to_str_wrong("hello".encode("utf-8"))
'hello'
\end{minted}

              \textbf{Example:} \mintinline{python}{b'\xe4\xb8\xad'}.\\
              \textbf{Answer:} This function is incorrect because it attempts to decode each byte individually, which fails for any character encoded as a multi-byte sequence, resulting in a \mintinline{python}{UnicodeDecodeError}.

        \item Give a two byte sequence that does not decode to any Unicode character(s).\\
              \textbf{Example:} \mintinline{python}{b'\xff\xff'}.\\
              \textbf{Answer:} This sequence is invalid because the byte \mintinline{python}{0xFF} is explicitly forbidden in the UTF-8 standard and does not correspond to any valid lead or continuation byte in the encoding scheme.

    \end{enumerate}

\end{problembox}
\newpage
\subsection{Subword Tokenization}

Sennrich et al. [2016] propose to use byte-pair encoding (BPE; Gage, 1994), a compression algorithm that iteratively replaces (``merges'') the most frequent pair of bytes with a single, new unused index.

\subsection{BPE Tokenizer Training}

\paragraph{Vocabulary initialization}

\paragraph{Pre-tokenization}
Use regular expressions to split the sentence into independent punctuation marks and independent items that may include spaces, and improve the merging effect by counting the occurrence frequency of these independent items.
\begin{minted}[xleftmargin=1.5em, framesep=5pt]{pycon}
>>> PAT = r"""'(?:[sdmt]|ll|ve|re)| ?\p{L}+| ?\p{N}+| ?[^\s\p{L}\p{N}]+|\s+(?!\S)|\s+"""
>>> # requires `regex` package
>>> import regex as re
>>> re.findall(PAT, "some text that i'll pre-tokenize")
['some', ' text', ' that', ' i', "'ll", ' pre', '-', 'tokenize']
\end{minted}

\paragraph{Compute BPE merges} During the BPE merging process, two \textbf{principles} must be followed:

\begin{itemize}[itemsep=0pt, parsep=0pt, topsep=0pt]
    \item Do not \textbf{cross} pre-tokenization boundaries (this is already guaranteed algorithmically);
    \item Handle issues with a \textbf{reproducible conflict resolution strategy}, such as selecting the lexicographically largest word pair to merge when the occurrence counts are the same.
\end{itemize}

\paragraph{Special tokens}

\newpage

\subsection{Experimenting with BPE Tokenizer Training}

\paragraph{Parallelizing pre-tokenization} Using the provided segmentation code and the \mintinline{python}{multiprocessing} library, parallelize the tokenization statistics.

\paragraph{Removing special tokens before pre-tokenization}

\paragraph{Optimizing the merging step} After performing a statistical pass to obtain a tokenization table with frequencies, perform merging only within that tokenization table.

\begin{tipbox}{Profiling}
    You should use profiling tools like \texttt{cProfile} or \texttt{scalene} to identify the bottlenecks in your implementation, and focus on optimizing those.
\end{tipbox}

\begin{problembox}{train\_bpe}{BPE Tokenizer Training (15 points)}

    \qquad \textbf{Deliverable:} Write a function that, given a path to an input text file, trains a (byte-level) BPE tokenizer. Your BPE training function should handle (at least) the following input parameters:

    \begin{itemize}[label={}, leftmargin=0pt, itemsep=0.8em, parsep=0pt]
        \item \texttt{input\_path}: \textcolor{green!50!black}{\texttt{str}} Path to a text file with BPE tokenizer training data.

        \item \texttt{vocab\_size}: \textcolor{green!50!black}{\texttt{int}} A positive integer that defines the maximum final vocabulary size (including the initial byte vocabulary, vocabulary items produced from merging, and any special tokens).

        \item \texttt{special\_tokens}: \textcolor{green!50!black}{\texttt{list[str]}} A list of strings to add to the vocabulary. These special tokens do not otherwise affect BPE training.
    \end{itemize}

    Your BPE training function should return the resulting vocabulary and merges:

    \begin{itemize}[label={}, leftmargin=0pt, itemsep=0.8em, parsep=0pt]
        \item \texttt{vocab}: \textcolor{green!50!black}{\texttt{dict[int, bytes]}} The tokenizer vocabulary, a mapping from \textcolor{green!50!black}{\texttt{int}} (token ID in the vocabulary) to \textcolor{green!50!black}{\texttt{bytes}} (token bytes).

        \item \texttt{merges}: \textcolor{green!50!black}{\texttt{list[tuple[bytes, bytes]]}} A list of BPE merges produced from training. Each list item is a \textcolor{green!50!black}{\texttt{tuple}} of \textcolor{green!50!black}{\texttt{bytes}} (\texttt{<token1>}, \texttt{<token2>}), representing that \texttt{<token1>} was merged with \texttt{<token2>}. The merges should be ordered by order of creation.
    \end{itemize}

\end{problembox}

\begin{ideabox}{BPE Tokenizer Training}

    \begin{itemize}[label={}, leftmargin=0pt, itemsep=0.8em, parsep=0pt]
        \item \textbf{Python plus Rust}: Python is used for rapid prototyping, while Rust handles the actual deployment of the product.

        \item \textbf{Glitch Tokens}: Due to insufficiently strict data cleaning, some meaningless words—namely \textbf{Mojibake}—may appear during the BPE training process. Moreover, because of BPE's \textbf{greedy, unsupervised, and non-semantic nature}, a large number of such meaningless word merges occur, turning them into abnormally long and nonsensical words.

        \item \textbf{Vocabulary Bloat (Dead Tokens)}: Because BPE is incremental and never removes merges, the intermediate products of long words permanently occupy the vocabulary, leading to wasted model parameters and tokenization fragmentation. In industry, there are three approaches to address this issue: First, absolutely aggressive data cleaning; second, stricter regex filtering. Both of these are common methods for limiting glitch tokens. The third approach is to switch to the \textbf{Unigram} algorithm, which is frequently used by Google-family models (such as T5 and some Llama variants).

        \item \textbf{Unigram}: Unigram works by "subtraction": it first initializes an extremely large vocabulary (e.g., 100,000 words) and \textbf{then calculates how much the overall corpus compression rate would be affected if a given word were removed}. Those useless intermediate composite words are directly eliminated and discarded during training.
    \end{itemize}

\end{ideabox}

\begin{problembox}{train\_bpe\_tinystories}{BPE Training on TinyStories (2 points)}

    \begin{enumerate}[label=(\alph*), itemsep=1em, topsep=0pt]

        \item Train a byte-level BPE tokenizer on the TinyStories dataset, using a maximum vocabulary size of 10,000. Make sure to add the TinyStories \texttt{\textless|endoftext|\textgreater} special token to the vocabulary. Serialize the resulting vocabulary and merges to disk for further inspection. How many hours and memory did training take? What is the longest token in the vocabulary? Does it make sense?

              \textbf{Resource requirements:} $\le$ 30 minutes (no GPUs), $\le$ 30GB RAM \\
              \textbf{Resource used:} 10.78 seconds, 2.6GB RAM \\
              \textbf{Longest token:} ' disappointment'

        \item Profile your code. What part of the tokenizer training process takes the most time?

              \textbf{Answer:} \mintinline{python}{process_chunk_with_special_tokens} takes up the vast majority of the time.

    \end{enumerate}

\end{problembox}

\begin{problembox}{train\_bpe\_expts\_owt}{BPE Training on OpenWebText (2 points)}

    \begin{enumerate}[label=(\alph*), itemsep=1em, topsep=0pt]

        \item Train a byte-level BPE tokenizer on the OpenWebText dataset, using a maximum vocabulary size of 32,000. Serialize the resulting vocabulary and merges to disk for further inspection. What is the longest token in the vocabulary? Does it make sense?

              \textbf{Resource requirements:} $\le$ 12 hours (no GPUs), $\le$ 100GB RAM \\
              \textbf{Resource used:} 6 minutes, 7.0GB RAM \\
              \textbf{Longest token:} '----------------------------------------------------------------'

        \item Compare and contrast the tokenizer that you get training on TinyStories versus OpenWebText.

              \textbf{Answer:} Because TinyStories has an extremely limited vocabulary, a 10k vocab size forces the BPE tokenizer to over-merge, creating tokens that represent entire multi-word phrases. In contrast, OpenWebText's diverse internet text requires the BPE algorithm to rely heavily on subwords (like prefixes and suffixes) to represent its massive vocabulary within a 32k limit, resulting in shorter, more fragmented tokens for rare words.

    \end{enumerate}

\end{problembox}

\subsection{BPE Tokenizer: Encoding and Decoding}

\subsubsection{Encoding text}

\paragraph{Special tokens.}

\paragraph{Memory considerations.} \vspace{-1em}

\subsubsection{Decoding text}

\begin{problembox}{tokenizer}{Implementing the tokenizer (15 points)}

    \qquad \textbf{Deliverable:} Implement a \texttt{Tokenizer} class that, given a vocabulary and a list of merges, encodes text into integer IDs and decodes integer IDs into text. Your tokenizer should also support user-provided special tokens (appending them to the vocabulary if they aren't already there).

    \qquad \textbf{Answer:} code in appendix.

\end{problembox}

\subsection{Experiments}

\begin{problembox}{tokenizer\_experiments}{Experiments with tokenizers \hfill (4 points)}

    \begin{enumerate}[label=(\alph*), itemsep=1em, topsep=0pt]

        \item Sample 10 documents from TinyStories and OpenWebText. Using your previously-trained TinyStories and OpenWebText tokenizers (10K and 32K vocabulary size, respectively), encode these sampled documents into integer IDs. What is each tokenizer's compression ratio (bytes/token)?

              \textbf{Answer:} The TinyStories tokenizer achieves a higher compression ratio (5.5 bytes/token) because its vocabulary easily memorizes long, frequent whole words from its restricted domain. In contrast, the OpenWebText tokenizer has a lower compression ratio (3.8 bytes/token) as it must rely heavily on shorter subwords to encode a highly diverse vocabulary.

        \item What happens if you tokenize your OpenWebText sample with the TinyStories tokenizer? Compare the compression ratio and/or qualitatively describe what happens.

              \textbf{Answer:} Tokenizing OpenWebText with the TinyStories tokenizer causes a severe domain mismatch, forcing the tokenizer to aggressively fragment unfamiliar web text into individual characters or very small byte sequences. Consequently, the token count explodes, resulting in a dramatically lower (worse) compression ratio compared to using the native OpenWebText tokenizer.

        \item Estimate the throughput of your tokenizer (e.g., in bytes/second). How long would it take to tokenize the Pile dataset (825GB of text)?

              \textbf{Answer:} The throughput of the tokenizer depends on the implementation details, but a rough estimate for a modern CPU might be around 100 MB/second. Given the Pile dataset's size of 825GB, it would take approximately 8.25 hours to tokenize the entire dataset.

        \item Using your TinyStories and OpenWebText tokenizers, encode the respective training and development datasets into a sequence of integer token IDs. We'll use this later to train our language model. We recommend serializing the token IDs as a NumPy array of datatype \texttt{uint16}. Why is \texttt{uint16} an appropriate choice?

              \textbf{Answer:} \texttt{uint16} is an appropriate choice because it can represent all possible token IDs within the range of 0 to 65,535, which is sufficient for most practical applications and allows for efficient memory usage.

    \end{enumerate}

\end{problembox}
\newpage

\section{Transformer Language Model Architecture}

\subsection{Transformer LM}

\subsubsection{Token Embeddings}

\subsubsection{Pre-norm Transformer Block}

\subsection{Output Normalization and Embedding}

\subsection{Remark: Batching, Einsum and Efficient Computation}

\subsubsection{Mathematical Notation and Memory Ordering}

\subsection{Basic Building Blocks: Linear and Embedding Modules}

\subsubsection{Parameter Initialization}

\subsubsection{Linear Module}

\begin{problembox}{linear}{Implementing the linear module \hfill (1 point)}

    \textbf{Deliverable:} Implement a \texttt{Linear} class that inherits from \texttt{torch.nn.Module} and performs a linear transformation. Your implementation should follow the interface of PyTorch's built-in \texttt{nn.Linear} module, except for not having a \texttt{bias} argument or parameter. We recommend the following interface:

    \begin{itemize}[label={}, leftmargin=0pt, itemsep=1em, parsep=0pt]

        \item \mintinline{python}{def __init__(self, in_features, out_features, device=None, dtype=None)} Construct a linear transformation module. This function should accept the following parameters:

              \begin{itemize}[label={}, leftmargin=1.5em, itemsep=0.5em, parsep=0pt, topsep=0.5em]
                  \item \texttt{in\_features}: \textcolor{green!50!black}{\texttt{int}} final dimension of the input
                  \item \texttt{out\_features}: \textcolor{green!50!black}{\texttt{int}} final dimension of the output
                  \item \texttt{device}: \textcolor{green!50!black}{\texttt{torch.device | None = None}} Device to store the parameters on
                  \item \texttt{dtype}: \textcolor{green!50!black}{\texttt{torch.dtype | None = None}} Data type of the parameters
              \end{itemize}

        \item \mintinline{python}{def forward(self, x: torch.Tensor) -> torch.Tensor} Apply the linear transformation to the input.

    \end{itemize}

    Make sure to:
    \begin{itemize}[itemsep=0.5em, parsep=0pt, topsep=0.5em]
        \item subclass \texttt{nn.Module}
        \item call the superclass constructor
        \item construct and store your parameter as $W$ (not $W^\top$) for memory ordering reasons, putting it in an \texttt{nn.Parameter}
        \item of course, don't use \texttt{nn.Linear} or \texttt{nn.functional.linear}
    \end{itemize}

    \vspace{0.5em}

    Code below;

\end{problembox}

\subsubsection{Embedding Module}

\begin{problembox}{embedding}{Implement the embedding module \hfill (1 point)}

    \textbf{Deliverable:} Implement the \texttt{Embedding} class that inherits from \texttt{torch.nn.Module} and performs an embedding lookup. Your implementation should follow the interface of PyTorch's built-in \texttt{nn.Embedding} module. We recommend the following interface:

    \begin{itemize}[label={}, leftmargin=0pt, itemsep=1em, parsep=0pt]

        \item \mintinline{python}{def __init__(self, num_embeddings, embedding_dim, device=None, dtype=None)} \\
              Construct an embedding module. This function should accept the following parameters:

              \begin{itemize}[label={}, leftmargin=1.5em, itemsep=0.5em, parsep=0pt, topsep=0.5em]
                  \item \texttt{num\_embeddings}: \textcolor{green!50!black}{\texttt{int}} Size of the vocabulary
                  \item \texttt{embedding\_dim}: \textcolor{green!50!black}{\texttt{int}} Dimension of the embedding vectors, i.e., $d_{\mathrm{model}}$
                  \item \texttt{device}: \textcolor{green!50!black}{\texttt{torch.device | None = None}} Device to store the parameters on
                  \item \texttt{dtype}: \textcolor{green!50!black}{\texttt{torch.dtype | None = None}} Data type of the parameters
              \end{itemize}

        \item \mintinline{python}{def forward(self, token_ids: torch.Tensor) -> torch.Tensor} \\
              Lookup the embedding vectors for the given token IDs.

    \end{itemize}

    Make sure to:
    \begin{itemize}[itemsep=0.5em, parsep=0pt, topsep=0.5em]
        \item subclass \texttt{nn.Module}
        \item call the superclass constructor
        \item initialize your embedding matrix as a \texttt{nn.Parameter}
        \item store the embedding matrix with the \texttt{d\_model} being the final dimension
        \item of course, don't use \texttt{nn.Embedding} or \texttt{nn.functional.embedding}
    \end{itemize}

    \vspace{0.5em}
    Again, use the settings from above for initialization, and use \texttt{torch.nn.init.trunc\_normal\_} to initialize the weights.

    To test your implementation, implement the test adapter at \textcolor{red!80!black}{\texttt{[adapters.run\_embedding]}}. Then, run \texttt{uv run pytest -k test\_embedding}.

\end{problembox}

\subsection{Pre-Norm Transformer Block}

\subsubsection{Root Mean Square Layer Normalization}

\begin{problembox}{rmsnorm}{Root Mean Square Layer Normalization (1 point)}

    \textbf{Deliverable:} Implement RMSNorm as a \texttt{torch.nn.Module}. We recommend the following interface:

    \begin{itemize}[label={}, leftmargin=0pt, itemsep=1em, parsep=0pt]

        \item \mintinline{python}{def __init__(self, d_model: int, eps: float = 1e-5, device=None, dtype=None)} \\
              Construct the RMSNorm module. This function should accept the following parameters:

              \begin{itemize}[label={}, leftmargin=1.5em, itemsep=0.5em, parsep=0pt, topsep=0.5em]
                  \item \texttt{d\_model}: \textcolor{green!50!black}{\texttt{int}} Hidden dimension of the model
                  \item \texttt{eps}: \textcolor{green!50!black}{\texttt{float = 1e-5}} Epsilon value for numerical stability
                  \item \texttt{device}: \textcolor{green!50!black}{\texttt{torch.device | None = None}} Device to store the parameters on
                  \item \texttt{dtype}: \textcolor{green!50!black}{\texttt{torch.dtype | None = None}} Data type of the parameters
              \end{itemize}

        \item \mintinline{python}{def forward(self, x: torch.Tensor) -> torch.Tensor} \\
              Process an input tensor of shape \texttt{(batch\_size, sequence\_length, d\_model)} and return a tensor of the same shape.

    \end{itemize}

    \vspace{0.5em}
    \textbf{Note:} Remember to upcast your input to \texttt{torch.float32} before performing the normalization (and later downcast to the original dtype), as described above.

    To test your implementation, implement the test adapter at \textcolor{red!80!black}{\texttt{[adapters.run\_rmsnorm]}}. Then, run \texttt{uv run pytest -k test\_rmsnorm}.

\end{problembox}

\subsubsection{Position-Wise Feed-Forward Network}

\begin{problembox}{positionwise\_feedforward}{Implement the position-wise feed-forward network (2 points)}

    \textbf{Deliverable:} \quad Implement the SwiGLU feed-forward network, composed of a SiLU activation function and a GLU.

    \vspace{0.5em}
    \textbf{Note:} in this particular case, you should feel free to use \texttt{torch.sigmoid} in your implementation for numerical stability.

    \vspace{0.5em}
    You should set $d_{\mathrm{ff}}$ to approximately $\frac{8}{3} \times d_{\mathrm{model}}$ in your implementation, while ensuring that the dimensionality of the inner feed-forward layer is a multiple of 64 to make good use of your hardware. To test your implementation against our provided tests, you will need to implement the test adapter at \textcolor{red!80!black}{\texttt{[adapters.run\_swiglu]}}. Then, run \texttt{uv run pytest -k test\_swiglu} to test your implementation.

\end{problembox}

\subsubsection{Relative Positional Embeddings}

\begin{problembox}{rope}{Implement RoPE (2 points)}

    \textbf{Deliverable:} Implement a class \texttt{RotaryPositionalEmbedding} that applies RoPE to the input tensor. The following interface is recommended:

    \begin{itemize}[label={}, leftmargin=0pt, itemsep=1em, parsep=0pt]

        \item \mintinline{python}{def __init__(self, theta: float, d_k: int, max_seq_len: int, device=None)} \\
              Construct the RoPE module and create buffers if needed.

              \begin{itemize}[label={}, leftmargin=1.5em, itemsep=0.5em, parsep=0pt, topsep=0.5em]
                  \item \texttt{theta}: \textcolor{green!50!black}{\texttt{float}} $\Theta$ value for the RoPE
                  \item \texttt{d\_k}: \textcolor{green!50!black}{\texttt{int}} dimension of query and key vectors
                  \item \texttt{max\_seq\_len}: \textcolor{green!50!black}{\texttt{int}} Maximum sequence length that will be inputted
                  \item \texttt{device}: \textcolor{green!50!black}{\texttt{torch.device | None = None}} Device to store the buffer on
              \end{itemize}

        \item \mintinline{python}{def forward(self, x: torch.Tensor, token_positions: torch.Tensor) -> torch.Tensor} \\
              Process an input tensor of shape \texttt{(..., seq\_len, d\_k)} and return a tensor of the same shape. Note that you should tolerate $x$ with an arbitrary number of batch dimensions. You should assume that the token positions are a tensor of shape \texttt{(..., seq\_len)} specifying the token positions of $x$ along the sequence dimension.

              \vspace{0.5em}
              You should use the token positions to slice your (possibly precomputed) \texttt{cos} and \texttt{sin} tensors along the sequence dimension.

    \end{itemize}

    \vspace{0.5em}
    To test your implementation, complete \textcolor{red!80!black}{\texttt{[adapters.run\_rope]}} and make sure it passes \texttt{uv run pytest -k test\_rope}.

\end{problembox}

\subsubsection{Scaled Dot-Product Attention}

\begin{problembox}{softmax}{Implement softmax \hfill (1 point)}

    \textbf{Deliverable:} Write a function to apply the softmax operation on a tensor. Your function should take two parameters: a tensor and a \textit{dimension} $i$, and apply softmax to the $i$-th dimension of the input tensor. The output tensor should have the same shape as the input tensor, but its $i$-th dimension will now have a normalized probability distribution. Use the trick of subtracting the maximum value in the $i$-th dimension from all elements of the $i$-th dimension to avoid numerical stability issues.

    To test your implementation, complete \textcolor{red!80!black}{\texttt{[adapters.run\_softmax]}} and make sure it passes \texttt{uv run pytest -k test\_softmax\_matches\_pytorch}.

\end{problembox}

\boxseparator

\begin{problembox}{scaled\_dot\_product\_attention}{Implement scaled dot-product attention \hfill (5 points)}

    \textbf{Deliverable:} \quad Implement the scaled dot-product attention function. Your implementation should handle keys and queries of shape \texttt{(batch\_size, \dots, seq\_len, d\_k)} and values of shape \texttt{(batch\_size, \dots, seq\_len, d\_v)}, where \dots represents any number of other batch-like dimensions (if provided). The implementation should return an output with the shape \texttt{(batch\_size, \dots, d\_v)}. See section 3.3 for a discussion on batch-like dimensions.

    Your implementation should also support an optional user-provided boolean mask of shape \texttt{(seq\_len, seq\_len)}. The attention probabilities of positions with a mask value of \texttt{True} should collectively sum to 1, and the attention probabilities of positions with a mask value of \texttt{False} should be zero.

    To test your implementation against our provided tests, you will need to implement the test adapter at \textcolor{red!80!black}{\texttt{[adapters.run\_scaled\_dot\_product\_attention]}}.

    \texttt{uv run pytest -k test\_scaled\_dot\_product\_attention} tests your implementation on third-order input tensors, while \texttt{uv run pytest -k test\_4d\_scaled\_dot\_product\_attention} tests your implementation on fourth-order input tensors.

\end{problembox}

\subsubsection{Causal Multi-Head Self-Attention}

\begin{problembox}{multihead\_self\_attention}{\newline Implement causal multi-head self-attention (5 points)}

    \textbf{Deliverable:} \quad Implement causal multi-head self-attention as a \texttt{torch.nn.Module}. Your implementation should accept (at least) the following parameters:

    \begin{itemize}[label={}, leftmargin=1.5em, itemsep=0.5em, parsep=0pt, topsep=0.5em]
        \item \texttt{d\_model}: \textcolor{green!50!black}{\texttt{int}} Dimensionality of the Transformer block inputs.
        \item \texttt{num\_heads}: \textcolor{green!50!black}{\texttt{int}} Number of heads to use in multi-head self-attention.
    \end{itemize}

    \vspace{0.5em}
    Following Vaswani et al. [2017], set $d_k = d_v = d_{\mathrm{model}}/h$. \newline To test your implementation against our provided tests, implement the test adapter at \textcolor{red!80!black}{\texttt{[adapters.run\_multihead\_self\_attention]}}. Then, run \texttt{uv run pytest -k test\_multihead\_self\_attention} to test your implementation.

\end{problembox}

\subsection{The Full Transformer LM}

\begin{problembox}{transformer\_block}{Implement the Transformer block (3 points)}

    Implement the pre-norm Transformer block as described in \S 3.5 and illustrated in Figure 2. Your Transformer block should accept (at least) the following parameters.

    \begin{itemize}[label={}, leftmargin=1.5em, itemsep=0.5em, parsep=0pt, topsep=0.5em]
        \item \texttt{d\_model}: \textcolor{green!50!black}{\texttt{int}} Dimensionality of the Transformer block inputs.
        \item \texttt{num\_heads}: \textcolor{green!50!black}{\texttt{int}} Number of heads to use in multi-head self-attention.
        \item \texttt{d\_ff}: \textcolor{green!50!black}{\texttt{int}} Dimensionality of the position-wise feed-forward inner layer.
    \end{itemize}

    \vspace{0.5em}
    To test your implementation, implement the adapter \textcolor{red!80!black}{\texttt{[adapters.run\_transformer\_block]}}. Then run \texttt{uv run pytest -k test\_transformer\_block} to test your implementation.

    \textbf{Deliverable:} Transformer block code that passes the provided tests.

\end{problembox}

\boxseparator

\begin{problembox}{transformer\_lm}{Implementing the Transformer LM  (3 points)}

    Time to put it all together! Implement the Transformer language model as described in \S 3.1 and illustrated in Figure 1. At minimum, your implementation should accept all the aforementioned construction parameters for the Transformer block, as well as these additional parameters:

    \begin{itemize}[label={}, leftmargin=1.5em, itemsep=0.5em, parsep=0pt, topsep=0.5em]
        \item \texttt{vocab\_size}: \textcolor{green!50!black}{\texttt{int}} The size of the vocabulary, necessary for determining the dimensionality of the token embedding matrix.
        \item \texttt{context\_length}: \textcolor{green!50!black}{\texttt{int}} The maximum context length, necessary for determining the dimensionality of the position embedding matrix.
        \item \texttt{num\_layers}: \textcolor{green!50!black}{\texttt{int}} The number of Transformer blocks to use.
    \end{itemize}

    \vspace{0.5em}
    To test your implementation against our provided tests, you will first need to implement the test adapter at \textcolor{red!80!black}{\texttt{[adapters.run\_transformer\_lm]}}. Then, run \texttt{uv run pytest -k test\_transformer\_lm} to test your implementation.

    \textbf{Deliverable:} A Transformer LM module that passes the above tests.

\end{problembox}

\boxseparator

\vspace{1em}
\noindent\textbf{Problem (transformer\_accounting): Transformer LM resource accounting \hfill (5 points)}

\noindent\rule{\textwidth}{0.4pt}

\questionitem{\textbf{Subproblem(a)}}
\begin{questionbox}
    Consider GPT-2 XL, which has the following configuration:

    \vspace{0.5em}
    \begin{tabular}{@{}ll}
        \texttt{vocab\_size}     & 50,257 \\
        \texttt{context\_length} & 1,024  \\
        \texttt{num\_layers}     & 48     \\
        \texttt{d\_model}        & 1,600  \\
        \texttt{num\_heads}      & 25     \\
        \texttt{d\_ff}           & 6,400
    \end{tabular}

    \vspace{0.5em}
    Suppose we constructed our model using this configuration. How many trainable parameters would our model have? Assuming each parameter is represented using single-precision floating point, how much memory is required to just load this model?
\end{questionbox}

% ================= 开始 Answerbox =================
\begin{answerbox}
    \textbf{Answer:}

    \textbf{A. Token Embeddings}
    \begin{itemize}[itemsep=0.5em, parsep=0pt, topsep=0.5em]
        \item \textbf{Calculation:} $\texttt{vocab\_size} \times d_{\mathrm{model}} = 50,257 \times 1,600 = \mathbf{80,411,200}$
        \item \textbf{Function:} Maps discrete word IDs to continuous numerical vectors; this is the first step in the model's understanding of language.
    \end{itemize}

    \vspace{1em}
    \textbf{B. Transformer Blocks (48 layers in total)}

    Each Block contains:
    \begin{enumerate}
        \item \textbf{Attention Projections:}
              \begin{itemize}[itemsep=0.2em, parsep=0pt, topsep=0.2em]
                  \item Four matrices $W_q, W_k, W_v, W_o$, each of size $d_{\mathrm{model}} \times d_{\mathrm{model}}$.
                  \item \textbf{Parameters:} $4 \times (1,600 \times 1,600) = 10,240,000$
              \end{itemize}

        \item \textbf{SwiGLU Feed-Forward Network (FFN):}
              \begin{itemize}[itemsep=0.2em, parsep=0pt, topsep=0.2em]
                  \item According to your implementation, it contains $W_1, W_3$ (up-projection) and $W_2$ (down-projection).
                  \item \textbf{Parameters:} $3 \times (d_{\mathrm{model}} \times d_{\mathrm{ff}}) = 3 \times (1,600 \times 6,400) = 30,720,000$
              \end{itemize}

        \item \textbf{Layer Normalization (RMSNorm):}
              \begin{itemize}[itemsep=0.2em, parsep=0pt, topsep=0.2em]
                  \item \texttt{ln1} and \texttt{ln2} each have a learnable scale parameter of length $d_{\mathrm{model}}$.
                  \item \textbf{Parameters:} $2 \times 1,600 = 3,200$
              \end{itemize}
    \end{enumerate}

    \begin{tabular}{@{}l@{}}
        \textbf{Subtotal per layer:} $10,240,000 + 30,720,000 + 3,200 = 40,963,200$ \\
        \textbf{Total for 48 layers:} $48 \times 40,963,200 = \mathbf{1,966,233,600}$
    \end{tabular}

    \vspace{1em}
    \textbf{C. Final Output Layer (Final Norm \& LM Head)}
    \begin{itemize}[itemsep=0.5em, parsep=0pt, topsep=0.5em]
        \item \textbf{Final RMSNorm:} $1,600$
        \item \textbf{LM Head:} Maps the hidden state back to the vocabulary, $d_{\mathrm{model}} \times \texttt{vocab\_size} = 1,600 \times 50,257 = 80,411,200$
        \item \textbf{Function:} Converts the high-dimensional features learned by the model back into a probability distribution over words, used to predict the next word.
    \end{itemize}

    \begin{tabular}{@{}l@{}}
        \textbf{Total Parameters:}                                         \\
        $80,411,200 + 1,966,233,600 + 80,411,200 = \mathbf{2,127,056,000}$ \\
        \textbf{Memory Cost(FP32):}                                        \\
        $2,127,057,600 \times 4 \text{ Bytes} \approx \mathbf{8,508,230,400} \text{ Bytes}$
    \end{tabular}

    \vspace{0.5em}
\end{answerbox}
% ================= 结束 Answerbox =================



\questionitem{\textbf{Subproblem(b)}}
\begin{questionbox}
    Identify the matrix multiplies required to complete a forward pass of our GPT-2 XL-shaped model. How many FLOPs do these matrix multiplies require in total? Assume that our input sequence has \texttt{context\_length} tokens.
\end{questionbox}

\begin{answerbox}
    \textbf{Answer:}

    \textbf{A. The Foyer (Entrance)}
    \begin{itemize}[itemsep=0.5em, parsep=0pt, topsep=0.5em]
        \item \textbf{Token Embedding:} Pure memory lookup to fetch the 1,600-dimensional vector. No arithmetic operations. \textbf{FLOPs: 0}.
        \item \textbf{Positional Encoding (RoPE):} Multiplies features by $\sin$ and $\cos$. Approx 4 operations per element.
              \newline \textbf{FLOPs:} $4 \times s \times d = 4 \times 1024 \times 1600 \approx \mathbf{6.55 \text{ MFLOPs}}$.
    \end{itemize}

    \vspace{1em}
    \textbf{B. Transformer Blocks (Single Layer)}

    Each Block contains:
    \begin{enumerate}
        \item \textbf{Attention Mechanism:}
              \begin{itemize}[itemsep=0.2em, parsep=0pt, topsep=0.2em]
                  \item \textbf{RMSNorm 1:} Approx 5 ops per element. \textbf{FLOPs:} $5 \times s \times d \approx 8 \text{ MFLOPs}$.
                  \item \textbf{Q, K, V Projections ($W_q, W_k, W_v$):} Complexity $\mathcal{O}(sd^2)$. \newline \textbf{FLOPs:} $3 \times 2sd^2 \approx \mathbf{15.73 \text{ GFLOPs}}$.
                  \item \textbf{Attention Scores \& Value Extraction:} $Q \times K^\top$ and $\text{Score} \times V$. \newline \textbf{FLOPs:} $2s^2d + 2s^2d \approx \mathbf{6.71 \text{ GFLOPs}}$.
                  \item \textbf{Internal Softmax:} Approx 3 ops per element. \textbf{FLOPs:} $3 \times s^2 \times h \approx 78.64 \text{ MFLOPs}$.
                  \item \textbf{Output Projection ($W_o$):} \textbf{FLOPs:} $2sd^2 \approx \mathbf{5.24 \text{ GFLOPs}}$.
                  \item \textbf{Residual Add 1:} \textbf{FLOPs:} $s \times d \approx 1.63 \text{ MFLOPs}$.
              \end{itemize}

        \item \textbf{SwiGLU Feed-Forward Network (FFN):}
              \begin{itemize}[itemsep=0.2em, parsep=0pt, topsep=0.2em]
                  \item \textbf{RMSNorm 2:} \textbf{FLOPs:} $\approx 8 \text{ MFLOPs}$.
                  \item \textbf{Up/Gate ($W_1, W_3$) \& Down ($W_2$) Projections:} \newline \textbf{FLOPs:} $(2 \times 2sdd_{\mathrm{ff}}) + (2sdd_{\mathrm{ff}}) = 6sdd_{\mathrm{ff}} \approx \mathbf{62.91 \text{ GFLOPs}}$.
                  \item \textbf{SiLU \& Gating:} Approx 5 ops per element. \textbf{FLOPs:} $5 \times s \times d_{\mathrm{ff}} \approx 32.76 \text{ MFLOPs}$.
                  \item \textbf{Residual Add 2:} \textbf{FLOPs:} $\approx 1.63 \text{ MFLOPs}$.
              \end{itemize}
    \end{enumerate}

    \begin{tabular}{@{}l@{}}
        \textbf{Single Block Summary:}                                        \\
        Large Matrix Multiplications: $\approx \mathbf{90.60 \text{ GFLOPs}}$ \\
        Minor Operations (Norms/Acts/Adds): $\approx \mathbf{137.21 \text{ MFLOPs}}$ (Only $\approx 0.15\%$ of a single layer)
    \end{tabular}

    \vspace{1em}
    \textbf{C. Scaling Up (48 Layers)}
    \begin{itemize}[itemsep=0.5em, parsep=0pt, topsep=0.5em]
        \item \textbf{Core Matrix Operations:} $48 \times 90.60 \text{ GFLOPs} \approx \mathbf{4,348.65 \text{ GFLOPs}}$.
        \item \textbf{Minor Operations:} $48 \times 137.21 \text{ MFLOPs} \approx \mathbf{6.58 \text{ GFLOPs}}$.
    \end{itemize}

    \vspace{1em}
    \textbf{D. The Exit (Final Output)}
    \begin{itemize}[itemsep=0.5em, parsep=0pt, topsep=0.5em]
        \item \textbf{Final RMSNorm:} \textbf{FLOPs:} $\approx 8 \text{ MFLOPs}$.
        \item \textbf{LM Head Projection:} Maps $1600$D features to vocabulary space. \newline \textbf{FLOPs:} $2 \times s \times d \times V = 2 \times 1024 \times 1600 \times 50,257 \approx \mathbf{164.68 \text{ GFLOPs}}$.
        \item \textbf{Final Softmax:} Approx 3 ops per element. \\
              \textbf{FLOPs:} $3 \times 1024 \times 50,257 \approx \mathbf{154.38 \text{ MFLOPs}}$.
    \end{itemize}

    \begin{tabular}{@{}l@{}}
        \textbf{Total FLOPs (Forward Pass):} \\
        $4348.65 + 6.58 + 164.68 + 0.15 \approx \mathbf{4,520 \text{ GFLOPs}} \text{ (approx. } \mathbf{4.52 \text{ TFLOPs}}\text{)}$
    \end{tabular}

    \vspace{0.5em}
    \rule{0pt}{1.5em}
\end{answerbox}



\questionitem{\textbf{Subproblem(c)}}
\begin{questionbox}
    Based on your analysis above, which parts of the model require the most FLOPs?
\end{questionbox}

\begin{answerbox}
    \textbf{Answer:}

    Based on the FLOPs analysis, the \textbf{Feed-Forward Network (FFN)} requires the absolute majority of the computation during a standard forward pass, rather than the Attention mechanism.

    \vspace{0.5em}
    \textbf{Key Insights:}
    \begin{enumerate}[itemsep=1em, parsep=0pt, topsep=0.5em]
        \item \textbf{The True Computational Bottleneck is the FFN, not Attention:}\\
              Due to the massive matrix dimension expansion (typically $d \rightarrow 4d$, or scaling to $d_{\mathrm{ff}}$ in SwiGLU), the computational cost of the FFN is nearly 3 times that of the attention linear projections ($Q, K, V, O$). This is precisely why cutting-edge large language models heavily research \textbf{Mixture of Experts (MoE)} architectures—sparsifying the FFN activations instantly saves a massive amount of compute!

        \item \textbf{The $\mathcal{O}(s^2)$ Complexity is Only Dominant at Long Contexts:}\\
              At a standard sequence length of $s=1024$, the quadratic attention operations (like $QK^\top$) actually account for a negligible fraction of the total computation. It is only when the sequence length $s$ scales up to tens or hundreds of thousands of tokens that this $\mathcal{O}(s^2)$ bottleneck overtakes the FFN to become the primary computational constraint.
    \end{enumerate}

    \vspace{0.5em}
    \rule{0pt}{1.5em}
\end{answerbox}



\questionitem{\textbf{Subproblem(d)}}
\begin{questionbox}
    Repeat your analysis with GPT-2 small (12 layers, 768 \texttt{d\_model}, 12 heads), GPT-2 medium (24 layers, 1024 \texttt{d\_model}, 16 heads), and GPT-2 large (36 layers, 1280 \texttt{d\_model}, 20 heads). As the model size increases, which parts of the Transformer LM take up proportionally more or less of the total FLOPs?

    \textbf{Deliverable:} For each model, provide a breakdown of model components and its associated FLOPs (as a proportion of the total FLOPs required for a forward pass). In addition, provide a one-to-two sentence description of how varying the model size changes the proportional FLOPs of each component.
\end{questionbox}

\begin{answerbox}
    \textbf{Answer:}

    Assuming the standard sequence length $s = 1024$, vocabulary size $V = 50,257$, and $d_{\mathrm{ff}} = 4 \times d_{\mathrm{model}}$, we calculate the core FLOPs (ignoring minor operations like Norms/Residuals which account for $<0.2\%$ of compute).

    \vspace{0.5em}
    \textbf{1. GPT-2 Small ($L=12, d_{\mathrm{model}}=768$)}
    \begin{itemize}[itemsep=0.2em, parsep=0pt, topsep=0.2em]
        \item \textbf{Total FLOPs:} $\approx 349.6 \text{ GFLOPs}$
        \item \textbf{Attention (Proj + Scores):} $\approx 96.7 \text{ GFLOPs}$ (\textbf{27.7\%})
        \item \textbf{FFN (SwiGLU):} $\approx 173.9 \text{ GFLOPs}$ (\textbf{49.7\%})
        \item \textbf{LM Head:} $\approx 79.0 \text{ GFLOPs}$ (\textbf{22.6\%})
    \end{itemize}

    \vspace{0.5em}
    \textbf{2. GPT-2 Medium ($L=24, d_{\mathrm{model}}=1024$)}
    \begin{itemize}[itemsep=0.2em, parsep=0pt, topsep=0.2em]
        \item \textbf{Total FLOPs:} $\approx 1,033.1 \text{ GFLOPs}$
        \item \textbf{Attention (Proj + Scores):} $\approx 309.2 \text{ GFLOPs}$ (\textbf{29.9\%})
        \item \textbf{FFN (SwiGLU):} $\approx 618.5 \text{ GFLOPs}$ (\textbf{59.9\%})
        \item \textbf{LM Head:} $\approx 105.4 \text{ GFLOPs}$ (\textbf{10.2\%})
    \end{itemize}

    \vspace{0.5em}
    \textbf{3. GPT-2 Large ($L=36, d_{\mathrm{model}}=1280$)}
    \begin{itemize}[itemsep=0.2em, parsep=0pt, topsep=0.2em]
        \item \textbf{Total FLOPs:} $\approx 2,257.8 \text{ GFLOPs}$
        \item \textbf{Attention (Proj + Scores):} $\approx 676.5 \text{ GFLOPs}$ (\textbf{30.0\%})
        \item \textbf{FFN (SwiGLU):} $\approx 1,449.6 \text{ GFLOPs}$ (\textbf{64.2\%})
        \item \textbf{LM Head:} $\approx 131.7 \text{ GFLOPs}$ (\textbf{5.8\%})
    \end{itemize}

    \vspace{0.5em}
    \textbf{Trend Description:} \\
    As the model size increases, the computational cost is heavily dominated by the internal Transformer blocks (specifically the FFN and Attention projections, which scale quadratically with the hidden dimension, $\mathcal{O}(L \cdot s \cdot d^2)$). Consequently, the LM Head (which scales only linearly with the hidden dimension, $\mathcal{O}(s \cdot d \cdot V)$) takes up a proportionally smaller fraction of the total FLOPs, plummeting from nearly $23\%$ in the Small model to roughly $6\%$ in the Large model.

    \vspace{0.5em}
    \rule{0pt}{1.5em}
\end{answerbox}



\questionitem{\textbf{Subproblem(e)}}
\begin{questionbox}
    Take GPT-2 XL and increase the context length to 16,384. How does the total FLOPs for one forward pass change? How do the relative contribution of FLOPs of the model components change?
\end{questionbox}

\begin{answerbox}
    \textbf{Answer:}

    Increasing the context length $s$ from $1,024$ to $16,384$ (a $16\times$ increase) causes the total FLOPs to skyrocket from roughly $4.5 \text{ TFLOPs}$ to \textbf{$\approx 149.5 \text{ TFLOPs}$} (an approximate $33\times$ increase).

    The relative contributions of the model components change fundamentally: the quadratic $\mathcal{O}(s^2d)$ attention score computation overtakes the FFN to become the absolute computational bottleneck.

    \vspace{0.5em}
    \textbf{Detailed Breakdown:}
    \begin{itemize}[itemsep=0.5em, parsep=0pt, topsep=0.5em]
        \item \textbf{Linear Scaling Components ($\mathcal{O}(s)$, scaled by $16\times$):}
              \begin{itemize}[itemsep=0.2em, parsep=0pt, topsep=0.2em]
                  \item \textbf{FFN:} $\approx 1006.6 \text{ GFLOPs}$ per layer $\times 48 \approx \mathbf{48.3 \text{ TFLOPs}}$ (\textbf{32.3\%})
                  \item \textbf{Linear Projections ($Q,K,V,O$):} $\approx 335.5 \text{ GFLOPs}$ per layer $\times 48 \approx \mathbf{16.1 \text{ TFLOPs}}$ (\textbf{10.8\%})
                  \item \textbf{LM Head:} $\approx \mathbf{2.6 \text{ TFLOPs}}$ (\textbf{1.8\%})
              \end{itemize}

        \item \textbf{Quadratic Scaling Component ($\mathcal{O}(s^2)$, scaled by $256\times$):}
              \begin{itemize}[itemsep=0.2em, parsep=0pt, topsep=0.2em]
                  \item \textbf{Attention Scores ($QK^\top$ \& $SV$):} The original $6.71 \text{ GFLOPs}$ per layer explodes to $\approx 1,718 \text{ GFLOPs}$ per layer! \\
                        Total across 48 layers: $\approx \mathbf{82.5 \text{ TFLOPs}}$ (\textbf{55.1\%})
              \end{itemize}
    \end{itemize}

    \vspace{0.5em}
    \textbf{Conclusion:}
    At $s=16,384$, the attention mechanism alone accounts for over \textbf{65\%} of the total compute (Projections + Scores), completely dwarfing the FFN. This demonstrates why techniques like FlashAttention, Ring Attention, or Sparse Attention are strictly required for long-context language models.

    \vspace{0.5em}
    \rule{0pt}{1.5em}
\end{answerbox}

\newpage

\section{Training a Transformer LM}

\subsection{Cross-entropy loss}

\noindent\textbf{Problem (cross\_entropy): Implement Cross entropy}

\noindent\rule{\textwidth}{0.4pt}

\begin{questionbox}
    \textbf{Deliverable:} Write a function to compute the cross entropy loss, which takes in predicted logits ($o_i$) and targets ($x_{i+1}$) and computes the cross entropy $\ell_i = -\log \mathrm{softmax}(o_i)[x_{i+1}]$. Your function should handle the following:

    \begin{itemize}[itemsep=0.5em, parsep=0pt, topsep=0.5em]
        \item Subtract the largest element for numerical stability.
        \item Cancel out log and exp whenever possible.
        \item Handle any additional batch dimensions and return the \textit{average} across the batch. As with section 3.3, we assume batch-like dimensions always come first, before the vocabulary size dimension.
    \end{itemize}

    \vspace{0.5em}
    Implement \texttt{[adapters.run\_cross\_entropy]}, then run \texttt{uv run pytest -k test\_cross\_entropy} to test your implementation.
\end{questionbox}

\subsection{The SGD Optimizer}

\subsubsection{Implementing SGD in PyTorch}

\noindent\textbf{Problem (learning\_rate\_tuning): Tuning the learning rate \hfill (1 point)}

\noindent\rule{\textwidth}{0.4pt}

\begin{questionbox}
    As we will see, one of the hyperparameters that affects training the most is the learning rate. Let's see that in practice in our toy example. Run the SGD example above with three other values for the learning rate: 1e1, 1e2, and 1e3, for just 10 training iterations. What happens with the loss for each of these learning rates? Does it decay faster, slower, or does it diverge (i.e., increase over the course of training)?

    \vspace{0.5em}
    \textbf{Deliverable:} A one-two sentence response with the behaviors you observed.
\end{questionbox}

\begin{answerbox}
    \textbf{Answer:}

    When the learning rates are \texttt{1e1} and \texttt{1e2}, the Loss decreases in both cases (with \texttt{1e1} exhibiting a smooth and rapid decay, while \texttt{1e2} shows a slight oscillation at the first step due to the larger step size before decaying rapidly). However, when the learning rate is \texttt{1e3}, the step size is too large, causing the update to directly overshoot the valley of the optimal solution. Consequently, the Loss increases rapidly during the training process and \textbf{diverges}.
\end{answerbox}

\boxseparator

\noindent\textbf{Problem (adamwAccounting): Resource accounting for training with AdamW \hfill (2 points)}

\noindent\rule{\textwidth}{0.4pt}

Let us compute how much memory and compute running AdamW requires. Assume we are using float32 for every tensor.

\questionitem{\textbf{Subproblem(a)}}
\begin{questionbox}
    How much peak memory does running AdamW require? Decompose your answer based on the memory usage of the parameters, activations, gradients, and optimizer state. Express your answer in terms of the \texttt{batch\_size} and the model hyperparameters (\texttt{vocab\_size}, \texttt{context\_length}, \texttt{num\_layers}, \texttt{d\_model}, \texttt{num\_heads}). Assume \texttt{d\_ff} $= 4 \times \texttt{d\_model}$.

    For simplicity, when calculating memory usage of activations, consider only the following components:
    \begin{itemize}
        \item Transformer block
              \begin{itemize}
                  \item RMSNorm(s)
                  \item Multi-head self-attention sublayer: $QKV$ projections, $Q^\top K$ matrix multiply, softmax, weighted sum of values, output projection.
                  \item Position-wise feed-forward: $W_1$ matrix multiply, SiLU, $W_2$ matrix multiply
              \end{itemize}
        \item final RMSNorm
        \item output embedding
        \item cross-entropy on logits
    \end{itemize}

    \textbf{Deliverable:} An algebraic expression for each of parameters, activations, gradients, and optimizer state, as well as the total.
\end{questionbox}

\begin{answerbox}
    \textbf{Answer:}

    To calculate the peak memory required for running AdamW, we must account for the static memory (model states) and dynamic memory (activations) saved during the forward pass for backpropagation. Let $P$ be the number of bytes per parameter (e.g., 2 or 4). We use the following abbreviations: $B$ (batch\_size), $S$ (context\_length), $V$ (vocab\_size), $L$ (num\_layers), $D$ (d\_model), and $H$ (num\_heads).

    \textbf{1. Static Memory (Model States)}
    \begin{itemize}
        \item \textbf{Parameters ($N$):} Token Embedding ($V \times D$) + $L$ Transformer Blocks ($2D$ for 2 RMSNorms + $3D^2$ for QKV + $D^2$ for Output Proj + $4D^2$ for FFN W1 + $4D^2$ for FFN W2) + Final RMSNorm ($D$) + Output Embedding ($V \times D$). \\
              Total $N = 2VD + L(12D^2 + 2D) + D$.
        \item \textbf{Gradients:} 1 memory slot per parameter $\rightarrow 1 \times N$.
        \item \textbf{AdamW Optimizer State:} Maintains 2 momentum states (first moment $m$ and second moment $v$) per parameter $\rightarrow 2 \times N$.
        \item \textbf{Total Static Memory:} $1N \text{ (Params)} + 1N \text{ (Grads)} + 2N \text{ (States)} = \mathbf{4N}$.
    \end{itemize}

    \textbf{2. Dynamic Memory (Activations, $A$)}
    \begin{itemize}
        \item \textbf{Transformer Block (per layer):}
              \begin{itemize}
                  \item RMSNorm (x2): Save inputs $\rightarrow 2 \times BSD$.
                  \item Self-Attention: QKV input ($BSD$) + $Q$ and $K$ for $Q^\top K$ multiply ($2BSD$) + Softmax output ($BHS^2$) + $V$ for weighted sum ($BSD$) + Output Proj input ($BSD$) $\rightarrow 5BSD + BHS^2$.
                  \item Position-wise FFN: W1 input ($BSD$) + pre-activation SiLU input ($4BSD$) + W2 input ($4BSD$) $\rightarrow 9BSD$.
              \end{itemize}
              \textit{Per Block Total:} $16BSD + BHS^2$.
        \item \textbf{Final Components:} Final RMSNorm input ($BSD$) + Output Embedding input ($BSD$) + Logits for Cross-Entropy ($BSV$).
        \item \textbf{Total Activations ($A$):} $\mathbf{L(16BSD + BHS^2) + 2BSD + BSV}$.
    \end{itemize}

    \textbf{3. Peak Memory Equation}
    Peak memory occurs at the end of the forward pass before gradients are freed, requiring both static states and activations to be stored simultaneously.
    $$ \text{Peak Memory (Bytes)} = P \times (4N + A) $$

    Expanded in terms of the hyperparameters:
    $$ \textbf{Peak Memory} = \mathbf{P \times \Big[ 4(2VD + 12LD^2 + 2LD + D) + L(16BSD + BHS^2) + 2BSD + BSV \Big]} $$

\end{answerbox}

\questionitem{\textbf{Subproblem(b)}}
\begin{questionbox}
    Instantiate your answer for a GPT-2 XL-shaped model to get an expression that only depends on the \texttt{batch\_size}. What is the maximum batch size you can use and still fit within 80GB memory?

    \textbf{Deliverable:} An expression that looks like $a \cdot \texttt{batch\_size} + b$ for numerical values $a, b$, and a number representing the maximum batch size.
\end{questionbox}

\begin{answerbox}
    \textbf{Answer:}

    We instantiate the memory formula for the GPT-2 XL model with the following hyperparameters: $L=48$, $D=1600$, $H=25$, $S=1024$, and $V=50257$. We assume $P=2$ bytes (FP16/BF16 training).

    \textbf{1. Static Memory Calculation:} \\
    First, calculate total parameters $N$:\\
    $N = 2VD + L(12D^2 + 2D) + D$\\
    $N = 2(50257)(1600) + 48(12 \cdot 1600^2 + 2 \cdot 1600) + 1600 \approx 1,635,536,000$ (1.64B parameters).\\
    Static Memory (Params + Grads + AdamW states) $= 4N \times P$ bytes:\\
    $4 \times 1.635 \times 10^9 \times 2 \approx 13,084,288,000 \text{ bytes} \approx \mathbf{13.08 \text{ GB}}$.\\

    \textbf{2. Activation Memory Calculation ($A$):} \\
    $A = B \times [L(16SD + HS^2) + 2SD + SV]$\\
    $16SD = 16(1024)(1600) = 26,214,400$\\
    $HS^2 = 25(1024^2) = 26,214,400$\\
    $SV = (1024)(50257) = 51,463,168$\\
    $A = B \times [48(26,214,400 + 26,214,400) + 3,276,800 + 51,463,168]$\\
    $A \approx B \times 2,571,322,368 \text{ elements}$.\\
    Activation Memory $= A \times P$ bytes:\\
    $2,571,322,368 \times 2 \approx 5,142,644,736 \text{ bytes} \times B \approx \mathbf{5.14 \text{ GB}} \cdot B$.\\

    \textbf{3. Final Expression and Max Batch Size:}
    The total peak memory expression is:
    $$ \text{Memory (GB)} = 5.14 \cdot \text{batch\_size} + 13.08 $$

    To fit within 80 GB:
    $5.14 \cdot B + 13.08 \le 80 \Rightarrow 5.14 \cdot B \le 66.92 \Rightarrow B \le 13.02$.
    The maximum batch size is \textbf{13}.

\end{answerbox}

\questionitem{\textbf{Subproblem(c)}}
\begin{questionbox}
    How many FLOPs does running one step of AdamW take?

    \textbf{Deliverable:} An algebraic expression, with a brief justification.
\end{questionbox}

\begin{answerbox}
    \textbf{Answer:}

    The total floating-point operations (FLOPs) for one step of AdamW training consist of the model computation (forward and backward passes) and the optimizer's parameter update.

    \textbf{1. Model Computation (Forward and Backward Passes):}
    For a Transformer model, the number of FLOPs is dominated by matrix multiplications.
    \begin{itemize}
        \item \textbf{Forward Pass:} Each parameter involves approximately one multiply-add operation per token, which equates to $2$ FLOPs. Total forward FLOPs $\approx 2NBS$.
        \item \textbf{Backward Pass:} The backward pass is approximately twice as computationally expensive as the forward pass because it must calculate gradients with respect to both the activations (to propagate the error) and the weights (to update the model). Total backward FLOPs $\approx 4NBS$.
    \end{itemize}
    Summing these gives the common heuristic of \textbf{$6NBS$} for the model computation.

    \textbf{2. AdamW Optimizer Update:}
    After the gradients are computed, the optimizer updates each of the $N$ parameters. Following the AdamW algorithm:
    \begin{itemize}
        \item First moment update ($m$): $2$ multiplications, $1$ addition = $3$ FLOPs.
        \item Second moment update ($v$): $1$ squaring, $2$ multiplications, $1$ addition = $4$ FLOPs.
        \item Parameter update: $1$ square root, $1$ division, $2$ multiplications, $2$ subtractions (including weight decay) $\approx 6$ FLOPs.
    \end{itemize}
    This adds a fixed cost of approximately \textbf{$13N$} FLOPs.

    \textbf{3. Final Expression:}
    $$ \mathbf{\textbf{Total FLOPs} \approx 6NBS + 13N} $$

    \textbf{Justification:}
    The majority of the workload ($6NBS$) comes from the tensor contractions during the forward and backward passes. While the AdamW update adds $13$ FLOPs per parameter, this term is negligible in large-scale training where the total number of tokens processed per step ($B \times S$) is much greater than $1$.

\end{answerbox}

\questionitem{\textbf{Subproblem(d)}}
\begin{questionbox}
    Model FLOPs utilization (MFU) is defined as the ratio of observed throughput (tokens per second) relative to the hardware's theoretical peak FLOP throughput [Chowdhery et al., 2022]. An NVIDIA A100 GPU has a theoretical peak of 19.5 teraFLOP/s for float32 operations. Assuming you are able to get 50\% MFU, how long would it take to train a GPT-2 XL for 400K steps and a batch size of 1024 on a single A100? Following Kaplan et al. [2020] and Hoffmann et al. [2022], assume that the backward pass has twice the FLOPs of the forward pass.

    \textbf{Deliverable:} The number of days training would take, with a brief justification.
\end{questionbox}

\begin{answerbox}
    \textbf{Answer:}

    To estimate the training duration, we first calculate the total computational work and then divide it by the effective hardware throughput.

    \textbf{1. Total Computational Work (FLOPs):}
    Using the heuristic that one training step (forward + backward pass) requires approximately $6N$ FLOPs per token:
    \begin{itemize}
        \item $N$ (GPT-2 XL parameters) $\approx 1.635 \times 10^9$
        \item $B$ (batch size) $= 1024$
        \item $S$ (context length) $= 1024$
        \item Total tokens per step $= B \times S = 1,048,576$
    \end{itemize}
    $\text{FLOPs per step} \approx 6 \times (1.635 \times 10^9) \times 1,048,576 \approx 1.028 \times 10^{16}$ FLOPs.
    For $T = 400,000$ steps, total work is:
    $\text{Total FLOPs} = 1.028 \times 10^{16} \times 400,000 \approx 4.11 \times 10^{21}$ FLOPs.

    \textbf{2. Effective Throughput:}
    The NVIDIA A100 peak is $19.5 \text{ TFLOPS} = 1.95 \times 10^{13}$ FLOPs/s. With an MFU of $50\%$:
    $\text{Effective Throughput} = 0.5 \times 1.95 \times 10^{13} = 9.75 \times 10^{12}$ FLOPs/s.

    \textbf{3. Training Time Calculation:}
    $\text{Time (seconds)} = \frac{4.11 \times 10^{21}}{9.75 \times 10^{12}} \approx 421,538,462$ seconds.
    Converting to days:
    $\text{Days} = \frac{421,538,462}{60 \times 60 \times 24} \approx \mathbf{4878.9 \text{ days}}$.

    \textbf{Justification:}
    Training a 1.6B parameter model for 400K steps involves massive compute ($4.11$ ZettaFLOPs). At 50\% utilization on a single A100, the hardware cannot process this volume in a reasonable timeframe, taking over 13 years. This underscores the necessity of multi-GPU distributed training for large language models.

\end{answerbox}

\subsection{Learning rate scheduling}

\noindent\textbf{Problem (learning\_rate\_schedule): Implement cosine learning rate schedule with warmup}

\noindent\rule{\textwidth}{0.4pt}

\begin{questionbox}
    Write a function that takes $t, \alpha_{\max}, \alpha_{\min}, T_w$ and $T_c$, and returns the learning rate $\alpha_t$ according to the scheduler defined above. Then implement \texttt{[adapters.get\_lr\_cosine\_schedule]} and make sure it passes \texttt{uv run pytest -k test\_get\_lr\_cosine\_schedule}.
\end{questionbox}

\subsection{Gradient clipping}

\noindent\textbf{Problem (gradient\_clipping): Implement gradient clipping \hfill (1 point)}

\noindent\rule{\textwidth}{0.4pt}

\begin{questionbox}
    Write a function that implements gradient clipping. Your function should take a list of parameters and a maximum $\ell_2$-norm. It should modify each parameter gradient in place. Use $\epsilon = 10^{-6}$ (the PyTorch default). Then, implement the adapter \texttt{[adapters.run\_gradient\_clipping]} and make sure it passes \texttt{uv run pytest -k test\_gradient\_clipping}.
\end{questionbox}


% \appendix
\section{Complete Source Code}

\subsection{Fast BPE Implementation (Rust)}
% 引入 Rust 代码
\inputminted[
    linenos,
    breaklines=true,
    fontsize=\footnotesize,
    frame=lines,
    framesep=2mm
]{rust}{../../fast_bpe/src/lib.rs}

\newpage

\subsection{Python BPE Trainer}
% 引入 Python 代码
\inputminted[
    linenos,               % 显示行号
    breaklines=true,       % 自动换行
    fontsize=\footnotesize,% 缩小字号，避免代码占用过多页数
    frame=lines,           % 上下加边框线
    framesep=2mm
]{python}{../test_function/bpe_trainer.py}

\newpage

\subsection{Tokenizer}
% 引入 Python 代码
\inputminted[
    linenos,               % 显示行号
    breaklines=true,       % 自动换行
    fontsize=\footnotesize,% 缩小字号，避免代码占用过多页数
    frame=lines,           % 上下加边框线
    framesep=2mm
]{python}{../test_function/tokenizer.py}

\newpage

\subsection{Adapters}
% 引入 Python 代码
\inputminted[
    linenos,               % 显示行号
    breaklines=true,       % 自动换行
    fontsize=\footnotesize,% 缩小字号，避免代码占用过多页数
    frame=lines,           % 上下加边框线
    framesep=2mm
]{python}{../../tests/adapters.py}

\end{document}
